%%%%%%%%%%%%%%%%
%%% gen 28 %%%
%%%%%%%%%%%%%%%%

\Chapter{28}

\Verse{1}
{
	\hP{וַיִּקְרָא יִצְחָק אֶל יַעֲקֹב וַיְבָרֶךְ אֹתוֹ וַיְצַוֵּהוּ וַיֹּאמֶר לוֹ לֹא תִקַּח אִשָּׁה מִבְּנוֹת כְּנָעַן׃}
}
{
	\eP{
		And Isaac called Jacob, and he blessed him and
		commanded him, and he said to him, “You shall not take a
		wife from the daughters of Canaan.
	}
}
{
	\aA{
		Ok, a quick recap before we proceed.
	}

	\aA{
		One sentence before this, Rebecca is disgusted with her life because of the daughters of Heth, which is the Hittites.
	}

	\aA{
		Two sentences before this, they're a bitterness of spirit to Rebecca.
	}

	\aA{
		Three sentences before this, we find out that Esau married a girl named Judith (``Jewish Girl'') and another one named Basemath (``Smells Good'') who are supposedly Hitties.
	}

	\aA{
		Now in this sentence, Isaac gives Jacob his blessing, and says he hopes that Jacob doesn't turn out how Esau did, specifically, he doesn't want him marrying any Canaanites, which ``includes'' the Hittites.\footnote{In the bible, Hittites are a type of Canaanites. That's an understandable inference, but not exactly true, as far as we can tell archaeologically. That is, since the text depicts Rebecca as not liking the Hittites, it's as if she's mad that her son married two girls named Sakura and Nahyun, and then the dad says to Jacob ``Don't go and marry a bunch of Chinese girls like your brother did.'' An understandable mistake. They're not exactly connoisseurs of the people in question.}
	}

	\aA{
		So Jacob gets his father's blessing fair and square because Jacob and Rebecca are mad that Esau married two birds from up north.
	}
}

\Verse{2}
{
	\hP{קוּם לֵךְ פַּדֶּנָה אֲרָם בֵּיתָה בְתוּאֵל אֲבִי אִמֶּךָ וְקַח לְךָ מִשָּׁם אִשָּׁה מִבְּנוֹת לָבָן אֲחִי אִמֶּךָ׃}
}
{
	\eP{
		Get up. Go to Paddan Aram, to the house of Bethuel, your
		mother’s father, and take a wife from there, from the
		daughters of Laban, your mother’s brother.
	}
}
{
	\aA{
		So that's why they tell Jacob to go to Laban's place and marry his cousins.
	}

	\aA{
		Recall that Laban means white.
	}

	\aA{
		So his daughters are white girls.
	}

	\aA{
		Also Laban lives Paddan Aram.
	}

	\aA{
		Remember that.
	}

	\aA{
		Jacob's going to Paddan Aram.
	}
}

\Verse{3}
{
	\hP{וְאֵל שַׁדַּי יְבָרֵךְ אֹתְךָ וְיַפְרְךָ וְיַרְבֶּךָ וְהָיִיתָ לִקְהַל עַמִּים׃}
}
{
	\eP{
		And may El Shadday bless you and make you fruitful and
		multiply you, so you’ll become a community of peoples,
	}
}
{
	\aA{
		And may God of the Mountain (or God of Breasts, we're not sure, but it's possible that Shaddai means something about fertility) give you lots of fertility (what a coincidence.)
	}
}

\Verse{4}
{
	\hP{וְיִתֶּן לְךָ אֶת בִּרְכַּת אַבְרָהָם לְךָ וּלְזַרְעֲךָ אִתָּךְ לְרִשְׁתְּךָ אֶת אֶרֶץ מְגֻרֶיךָ אֲשֶׁר נָתַן אֱלֹהִים לְאַבְרָהָם׃}
}
{
	\eP{
		and may He give you the blessing of Abraham, to you and to
		your seed with you, for you to possess the land of your
		residences, which {\God} gave to Abraham.”
	}
}
{
	\aA{
		He goes on about seed for awhile.
	}
}

\Verse{5}
{
	\hP{וַיִּשְׁלַח יִצְחָק אֶת יַעֲקֹב וַיֵּלֶךְ פַּדֶּנָה אֲרָם אֶל לָבָן בֶּן בְּתוּאֵל הָאֲרַמִּי אֲחִי רִבְקָה אֵם יַעֲקֹב וְעֵשָׂו׃}
}
{
	\eP{
		And Isaac sent Jacob, and he went to Paddan Aram, to
		Laban, son of Bethuel, the Aramean, brother of Rebekah,
		mother of Jacob and Esau.
	}
}
{
	\aA{
		Jacob goes to the White house.
	}

	\aA{
		It's in Paddan Aram.
	}
}

\Verse{6}
{
	\hP{וַיַּרְא עֵשָׂו כִּי בֵרַךְ יִצְחָק אֶת יַעֲקֹב וְשִׁלַּח אֹתוֹ פַּדֶּנָה אֲרָם לָקַחַת לוֹ מִשָּׁם אִשָּׁה בְּבָרְכוֹ אֹתוֹ וַיְצַו עָלָיו לֵאמֹר לֹא תִקַּח אִשָּׁה מִבְּנוֹת כְּנָעַן׃}
}
{
	\eP{
		And Esau saw that Isaac had blessed Jacob and sent him to
		Paddan Aram to take a wife for himself from there when he
		blessed him, and he had commanded him, saying, “You shall
		not take a wife from the daughters of Canaan,”
	}
}
{
%	\footnote{
%		sent. Reading the verb as Hebrew Qal (“and he sent him”) as in v. 5,
%		rather than as Piel (“and he let him go”) with the MT.
%	}

	\aA{
		Esau sees that Jacob got the blessing because of the whole Hittite Wives thing.
	}
}

\Verse{7}
{
	\hP{וַיִּשְׁמַע יַעֲקֹב אֶל אָבִיו וְאֶל אִמּוֹ וַיֵּלֶךְ פַּדֶּנָה אֲרָם׃}
}
{
	\eP{
		and Jacob had listened to his father and to his mother and
		had gone to Paddan Aram.
	}
}
{
	\aA{
		Jacob had already gone to the White House.
	}

	\aA{
		It's in Paddan Aram.
	}
}

\Verse{8}
{
	\hP{וַיַּרְא עֵשָׂו כִּי רָעוֹת בְּנוֹת כְּנָעַן בְּעֵינֵי יִצְחָק אָבִיו׃}
}
{
	\eP{
		And Esau saw that the daughters of Canaan were bad in his
		father Isaac’s eyes.
	}
}
{
	\aA{
		Esau sees that his dad doesn't like Canaanite girls.
	}
}

\Verse{9}
{
	\hP{וַיֵּלֶךְ עֵשָׂו אֶל יִשְׁמָעֵאל וַיִּקַּח אֶת מָחֲלַת בַּת יִשְׁמָעֵאל בֶּן אַבְרָהָם אֲחוֹת נְבָיוֹת עַל נָשָׁיו לוֹ לְאִשָּׁה׃}
}
{
	\eP{
		And Esau went to Ishmael and took Mahalath, daughter of
		Ishmael, son of Abraham, sister of Nebaioth, in addition
		to his wives as a wife for him.
	}
}
{
	\aA{
		So Esau goes to his uncle and marries his cousin as a third wife, which he hopes his dad will approve of.
	}
}

\Verse{10}
{
	\hJ{וַיֵּצֵא יַעֲקֹב מִבְּאֵר שָׁבַע וַיֵּלֶךְ חָרָנָה׃}
}
{
	\eJ{And Jacob left Beer-sheba and went to Haran.}
}
{
	\aA{
		Okay so at this point Jacob leaves.
	}

	\aA{
		He goes to Haran.
	}
}

\Verse{11}
{
	\hJ{וַיִּפְגַּע בַּמָּקוֹם וַיָּלֶן שָׁם כִּי בָא}
	\hE{הַשֶּׁמֶשׁ וַיִּקַּח מֵאַבְנֵי הַמָּקוֹם וַיָּשֶׂם מְרַאֲשֹׁתָיו וַיִּשְׁכַּב בַּמָּקוֹם הַהוּא׃}
}
{
	\eJ{And he happened upon a place and stayed the night there because}
	\eE{the sun was setting. And he took from the stones of the place
		and set it as his headrest and lay down in that place.}
}
{
	\aA{
		It's bedtime.
		
		So he stays the night.

		He uses a rock for a pillow.
	}
}

\Verse{12}
{
	\hE{וַיַּחֲלֹם וְהִנֵּה סֻלָּם מֻצָּב אַרְצָה וְרֹאשׁוֹ מַגִּיעַ הַשָּׁמָיְמָה וְהִנֵּה מַלְאֲכֵי אֱלֹהִים עֹלִים וְיֹרְדִים בּוֹ׃}
}
{
	\eE{
		And he dreamed. And here was a ladder, set up on the
		earth, and its top reaching to the skies. And here were
		angels of {\God}, going up and going down by it.
	}
}
{
	\aA{
		He dreams about a ladder to heaven.
	}

	\aA{
		There's angels of God going up and down it.
	}
}

\Verse{13}
{
	\hJ{וְהִנֵּה יְהֹוָה נִצָּב עָלָיו וַיֹּאמַר אֲנִי יְהֹוָה אֱלֹהֵי אַבְרָהָם אָבִיךָ וֵאלֹהֵי יִצְחָק הָאָרֶץ אֲשֶׁר אַתָּה שֹׁכֵב עָלֶיהָ לְךָ אֶתְּנֶנָּה וּלְזַרְעֶךָ׃}
}
{
	\eJ{
		And here was YHWH standing over him, and He said, “I am
		YHWH, your father Abraham’s {\God} and Isaac’s {\God}. The land
		on which you’re lying: I’ll give it to you and to your
		seed.
	}
}
{
	\aA{
		Okay so YHWH is standing over him. He does The seed schtick again.
	}
}

\Verse{14}
{
	\hJ{וְהָיָה זַרְעֲךָ כַּעֲפַר הָאָרֶץ וּפָרַצְתָּ יָמָּה וָקֵדְמָה וְצָפֹנָה וָנֶגְבָּה וְנִבְרְכוּ בְךָ כּל מִשְׁפְּחֹת הָאֲדָמָה וּבְזַרְעֶךָ׃}
}
{
	\eJ{
		And your seed will be like the dust of the earth, and
		you’ll expand to the west and east and north and south,
		and all the families on the earth will be blessed through
		you and through your seed.
	}
}
{
%	\footnote{
%		all the families on the earth. The prediction that God made to Abraham
%		and Isaac that all humankind will be blessed through them is now
%		repeated to Jacob as well. The confirmation of this prediction to each
%		of the three patriarchs individually highlights its tremendous
%		importance.
%	}

	\aA{
		He's like Jacob, look at everything.

		Everything you just seed will be yours.

		He goes on for a while about the seed.
	}
}

\Verse{15}
{
	\hJ{וְהִנֵּה אָנֹכִי עִמָּךְ וּשְׁמַרְתִּיךָ בְּכֹל אֲשֶׁר תֵּלֵךְ וַהֲשִׁבֹתִיךָ אֶל הָאֲדָמָה הַזֹּאת כִּי לֹא אֶעֱזבְךָ עַד אֲשֶׁר אִם עָשִׂיתִי אֵת אֲשֶׁר דִּבַּרְתִּי לָךְ׃}
}
{
	\eJ{
		And here I am with you, and I’ll watch over you everywhere
		that you’ll go, and I’ll bring you back to this land, for
		I won’t leave you until I’ve done what I’ve spoken to
		you.”
	}
}
{
	\aA{
		He says he's never going to give him up, never going to let him down, never going to run around or desert him.
	}
}

\Verse{16}
{
	\hJ{וַיִּיקַץ יַעֲקֹב מִשְּׁנָתוֹ וַיֹּאמֶר אָכֵן יֵשׁ יְהֹוָה בַּמָּקוֹם הַזֶּה וְאָנֹכִי לֹא יָדָעְתִּי׃}
}
{
	\eJ{
		And Jacob woke from his sleep and said, “YHWH is actually
		in this place, and I didn’t know!”
	}
}
{
	\aA{
		Jacob woke with a start.
	}

	\aA{
		YHWH's here.
	}

}

\Verse{17}
{
	\hE{וַיִּירָא וַיֹּאמַר מַה נּוֹרָא הַמָּקוֹם הַזֶּה אֵין זֶה כִּי אִם בֵּית אֱלֹהִים וְזֶה שַׁעַר הַשָּׁמָיִם׃}
}
{
	\eE{
		And he was afraid, and he said, “How awesome this place
		is! This is none other than {\God}’s house, and this is the
		gate of the skies!”
	}
}
{
	\aA{
		So he's afraid.
	}

	\aA{
		He's like ``This place is awesome!''
	}

	\aA{
		It's basically Beth-el.

		Like Heaven's Gate.
	}
}

\Verse{18}
{
	\hE{וַיַּשְׁכֵּם יַעֲקֹב בַּבֹּקֶר וַיִּקַּח אֶת הָאֶבֶן אֲשֶׁר שָׂם מְרַאֲשֹׁתָיו וַיָּשֶׂם אֹתָהּ מַצֵּבָה וַיִּצֹק שֶׁמֶן עַל רֹאשָׁהּ׃}
}
{
	\eE{
		And Jacob got up early in the morning and took the stone
		that he had set as his headrest and set it as a pillar and
		poured oil on its top.
	}
}
{
	\aA{
		Jacob turns his rock pillow up vertical and pours oil on it.
	}

	\aA{
		It's not clear why.%
		\aC{\footnote{Jacob is turning his pillow-rock into a sacred \textbf{standing stone} (\emph{maṣṣēbâ}): he stands it upright and pours oil over it to consecrate it. This was a normal ancient West Semitic way of marking a divine presence or holy place; Jacob does it again at Bethel (Gen 35:14), and erects other standing stones in Gen 31:45 and 35:20. The point becomes explicit in 28:22: “this stone … shall be God’s house”—part of the story’s explanation of \textbf{Bethel}, \emph{bêt-ʾēl}, “House of God/El.” Strikingly, Deut 16:22 later outright bans erecting a \emph{maṣṣēbâ}.}}
	}
}

\Verse{19}
{
	\hJ{וַיִּקְרָא אֶת שֵׁם הַמָּקוֹם הַהוּא בֵּית אֵל וְאוּלָם לוּז שֵׁם הָעִיר לָרִאשֹׁנָה׃}
}
{
	\eJ{
		And he called that place’s name Beth-El, though in fact
		Luz was the name of the city at first.
	}
}
{
	\aA{
		Oh wow it is Beth-el.
	}

	\aA{
		I was just joking earlier.
	}

	\aA{
		Noted.
	}
}

\Verse{20}
{
	\hE{וַיִּדַּר יַעֲקֹב נֶדֶר לֵאמֹר אִם יִהְיֶה אֱלֹהִים עִמָּדִי וּשְׁמָרַנִי בַּדֶּרֶךְ הַזֶּה אֲשֶׁר אָנֹכִי הוֹלֵךְ וְנָתַן לִי לֶחֶם לֶאֱכֹל וּבֶגֶד לִלְבֹּשׁ׃}
}
{
	\eE{
		And Jacob made a vow, saying, “If {\God} will be with me and
		watch over me in this way that I’m going and give me bread
		to eat and clothing to wear,
	}
}
{
	\aA{
		He's like ``If God will be with me, take care of me, give me bread, and give me clothes,''
	}
}

\Verse{21}
{
	\hE{וְשַׁבְתִּי בְשָׁלוֹם אֶל בֵּית אָבִי וְהָיָה יְהֹוָה לִי לֵאלֹהִים׃}
}
{
	\eE{
		and I come back in peace to my father’s house, then YHWH
		will become my {\God},
	}
}
{
	\aA{
		``And keep me safe until I get back to daddy, then YHWH will become my God,''\aC{\footnote{The 5th out of 5 uses of the name YHWH outside of J prior to the big reveal of the name of god in Exodus 3 (and 6). The Evidence: Dreams are heavily E-associated: Gen 20:3, 6; 28:12; 31:10–11, 24; 37:5–10; 40:5–23; 41:1–32. Standing stones also recur in E-linked Jacob material: Gen 28:18, 22; 31:45, 51–52; 35:14. And the actual sentence says, “If Elohim is with me ... then YHWH will be my God.” Jacob does not start by casually calling Elohim “YHWH”; YHWH appears in the result of the condition. E may therefore be deliberately saying: if the Elohim who appeared to Jacob actually protects him, Jacob will accept YHWH as his God.}}
	}
}

\Verse{22}
{
	\hE{וְהָאֶבֶן הַזֹּאת אֲשֶׁר שַׂמְתִּי מַצֵּבָה יִהְיֶה בֵּית אֱלֹהִים וְכֹל אֲשֶׁר תִּתֶּן לִי עַשֵּׂר אֲעַשְּׂרֶנּוּ לָךְ׃}
}
{
	\eE{
		and this stone that I set as a pillar will be {\God}’s house,
		and everything that you’ll give me I’ll tithe to you.”
	}
}
{
	\aA{
		``and my rock pillow will be God's house, and everything you give me, I'll give you 10\%.''
	}

	\aA{
		End of chapter.
	}
}
